%=============================================================================!
% 8.4  FORCES, THE VIRIAL, AND CHECKING FORCES                                !
%=============================================================================!
\section{Forces, the virial, and checking forces}
\label{sec:pe-forces}

A simulation needs the force on every atom at every step, and a few other
sums computed alongside it. This section derives them for a pair
potential, shows the code that computes them, and shows how to check any
force routine against its energy.

\subsection*{Forces}

Atom $i$ appears in the sum \cref{eq:pe-pair-sum} in the $N - 1$ terms
$\varphi(r_{ij})$ with $j \ne i$. \Cref{sec:ro-system} found the gradient
of one such term, $-\grad_i\varphi(r_{ij}) = \varphi'(r_{ij})\,\vb{r}_{ij}/
r_{ij}$ with $\vb{r}_{ij} = \vb{r}_j - \vb{r}_i$, so
\begin{equation}
   \vb{F}_i = \sum_{j\ne i}\varphi'(r_{ij})\,\frac{\vb{r}_{ij}}{r_{ij}} .
   \label{eq:pe-pair-force}
\end{equation}
Each term is the force on $i$ from $j$, along the line joining them,
towards $j$ when $\varphi' > 0$, where the energy rises with distance and
the pair attracts. The force on $j$ from $i$ is the same with $i$ and $j$
exchanged, and since $\vb{r}_{ji} = -\vb{r}_{ij}$ it is equal and
opposite: Newton's third law (\cref{sec:ne-third}). A program therefore
computes each pair once and adds its force to one atom and subtracts it
from the other, which is what \texttt{pair\_energy\_forces} of
\texttt{mdlab} does:

\begin{code}
def pair_energy_forces(
    positions: ArrayLike, pair: PairFunction
) -> tuple[float, NDArray, float]:
    r = np.asarray(positions, dtype=float)
    n = len(r)
    i, j = np.triu_indices(n, k=1)  # every pair once
    rij = r[i] - r[j]
    dist = np.linalg.norm(rij, axis=1)
    phi, dphi = pair(dist)
    f_on_i = -(dphi / dist)[:, None] * rij  # force on i from j
    forces = np.zeros_like(r)
    np.add.at(forces, i, f_on_i)
    np.add.at(forces, j, -f_on_i)
    return float(phi.sum()), forces, float(-(dist * dphi).sum())
\end{code}

\noindent \texttt{np.triu\_indices(n, k=1)} lists the pairs $(i, j)$ with
$i < j$ as two arrays of the same length; \texttt{rij} holds $\vb{r}_i -
\vb{r}_j = -\vb{r}_{ij}$ for every pair at once, so the minus sign makes
\texttt{f\_on\_i} the force of \cref{eq:pe-pair-force}; and
\texttt{np.add.at} adds each pair's force into the row of its atom, even
when an atom appears in many pairs; the plainer \texttt{forces[i] +=
f\_on\_i} would keep only one force for an atom listed several times in
\texttt{i}. The argument \texttt{pair} is any
function that returns $\varphi$ and $\varphi'$ for an array of distances,
such as \texttt{lennard\_jones}.

\subsection*{The virial}

The last number returned is the virial, $\mathcal{W} = \sum_i\vb{r}_i\cdot
\vb{F}_i$, from which Chapter~14 computes the pressure. For a pair potential it
takes a simple form. The pair $(i, j)$ contributes $\vb{F}_{ij}$ to atom
$i$ and $-\vb{F}_{ij}$ to atom $j$, and so $\vb{r}_i\cdot\vb{F}_{ij} -
\vb{r}_j\cdot\vb{F}_{ij} = -\vb{r}_{ij}\cdot\vb{F}_{ij}$ to $\mathcal{W}$. With
$\vb{F}_{ij} = \varphi'(r_{ij})\,\vb{r}_{ij}/r_{ij}$ and $\vb{r}_{ij}\cdot
\vb{r}_{ij} = r_{ij}^2$,
\begin{equation}
   \mathcal{W} = -\sum_{i<j} r_{ij}\,\varphi'(r_{ij}) .
   \label{eq:pe-virial}
\end{equation}
It depends only on distances, so it does not change when the origin is
moved. That holds for any potential whose forces add to zero: moving the
origin by $\vb{a}$ replaces each $\vb{r}_i$ by $\vb{r}_i - \vb{a}$ and
changes $\mathcal{W}$ by $-\vb{a}\cdot\sum_i\vb{F}_i = 0$ (\cref{sec:ha-noether}).
A pair pushed closer than its minimum repels, $\varphi' < 0$, and
contributes positively; a stretched pair attracts and contributes
negatively.

\subsection*{Checking forces}

A force routine is easy to get wrong and hard to check by eye, and a force
that is not exactly $-\grad U$ makes the energy drift (\cref{sec:en-drift}).
The standard check compares it with central differences of the energy,
$F \approx -[U(x + h) - U(x - h)]/(2h)$ for each coordinate in turn
(\cref{sec:mo-atoms}). Two errors limit the check. The Taylor series of
\cref{sec:mo-taylor} gives $U(x \pm h) = U \pm hU' + \frac12 h^2U'' \pm
\frac16 h^3U''' + \frac1{24}h^4U'''' + \order{h^5}$; subtracting cancels
the even powers, which have the same sign in both, as in
\cref{sec:mo-atoms}, so
\begin{equation*}
   \frac{U(x + h) - U(x - h)}{2h} = U' + \frac{h^2}{6}\,U''' + \order{h^4} ,
\end{equation*}
an error that shrinks as $h^2$. But a computer stores a number with about
16 significant decimal digits: rounding to the nearest stored number
changes it by at most half the machine epsilon, $\epsilon_{\mathrm m} =
2.2\times10^{-16}$ (the gap between 1 and the next stored number), times
its size, so a computed energy, a sum of many rounded terms, is uncertain
by about $\epsilon_{\mathrm m}\lvert U\rvert$ or more. The difference of two such energies
divided by $2h$ carries an error of order $\epsilon_{\mathrm m}\lvert U
\rvert/h$, which grows as $h$ shrinks. The total,
\begin{equation}
   \text{error} \approx \frac{h^2}{6}\,\lvert U'''\rvert
   + \frac{\epsilon_{\mathrm m}\lvert U\rvert}{h} ,
   \label{eq:pe-fd-error}
\end{equation}
is least where its slope with respect to $h$, $h\lvert U'''\rvert/3 -
\epsilon_{\mathrm m}\lvert U\rvert/h^2$, is zero, that is at $h^3 =
3\epsilon_{\mathrm m}\lvert U\rvert/\lvert U'''\rvert$: a step of about
the cube root of $\epsilon_{\mathrm m}$, $\num{6e-6}$, times the length on
which $U$ changes (if $U$ changes by about its own size over a length
$L$, then $\lvert U'''\rvert \approx \lvert U\rvert/L^3$ and $h \approx
(3\epsilon_{\mathrm m})^{1/3}L$); near the steep repulsion $L$ is much
shorter than $\sigma$ (\cref{ex:pe-best-step}). \Cref{fig:pe-fd} measures it for twelve Lennard-Jones
atoms. The error falls as $h^2$ (fitted slope 2.03) for large steps,
grows as $1/h$ (slope $-1.01$) for small ones, and is least,
$\num{1.0e-8}\,\varepsilon/\sigma$, at $h = 10^{-6}\sigma$, against forces
of up to $129\,\varepsilon/\sigma$.

\begin{figure}[htb]
\centering
\includegraphics{ch08_potentials/fd}
\caption{The largest error of forces from central differences of the
energy, for twelve Lennard-Jones atoms, against the step $h$: falling as
$h^2$ (dashed) from the Taylor series and rising as $1/h$ (dashed) from
the rounding of the energies.}
\label{fig:pe-fd}
\end{figure}

\begin{practice}
Every new force routine, and every change to one, should be checked
against central differences of its own energy, on an arrangement of atoms
that is not symmetric, so that no error cancels. A step near
$10^{-6}$ times the distance between atoms, the best step of
\cref{fig:pe-fd}, is a good start; the check should be repeated at a step
ten times larger and smaller, and the difference should behave as
\cref{eq:pe-fd-error} predicts. The tests of \texttt{mdlab} compare every
force routine of this chapter with central differences, at one step.
\end{practice}

\begin{keyidea}
For a pair potential, $\vb{F}_i = \sum_j\varphi'(r_{ij})\,\vb{r}_{ij}/
r_{ij}$, computed once per pair and applied with opposite signs, and the
virial is $\mathcal{W} = -\sum r_{ij}\varphi'(r_{ij})$. Forces are checked against
central differences of the energy, whose error is least at a step near
$\epsilon_{\mathrm m}^{1/3}$ times the length scale.
\end{keyidea}

\notebook{08}{break a force routine on purpose and catch it}

\begin{selfcheck}
Two Lennard-Jones atoms sit at $r = \sigma$. Is their virial positive or
negative?
\end{selfcheck}
